% Generated draft: requires mathematical and source review.
\documentclass[11pt]{article}
\usepackage[margin=25mm]{geometry}
\usepackage{fontspec}
\setmainfont{Latin Modern Roman}
\usepackage{amsmath,amssymb,mathtools,unicode-math}
\setmathfont{Latin Modern Math}
\usepackage{enumitem,needspace,xurl,hyperref}
\setlength{\parindent}{0pt}
\setlength{\parskip}{5pt}
\setlength{\emergencystretch}{3em}
\newcommand{\N}{\mathbb{N}}
\newcommand{\Z}{\mathbb{Z}}
\newcommand{\Q}{\mathbb{Q}}
\newcommand{\R}{\mathbb{R}}
\newcommand{\C}{\mathbb{C}}
\newcommand{\F}{\mathbb{F}}
\newcommand{\A}{\mathbb{A}}
\newcommand{\PP}{\mathbb{P}}
\newcommand{\E}{\mathbb{E}}
\newcommand{\eps}{\varepsilon}
\newcommand{\dd}{\,\mathrm{d}}
\begin{document}
\section*{UnsolvedMath ID: 30002070}
\textbf{Source problem code:} catalogue entry
\subsection{Definitions and mathematical statement}
A simplicial triangulation of a topological \(3\)-manifold \(M\) is a simplicial complex \(K\) together with a homeomorphism \(|K|\cong M\). For an edge \(e\) of \(K\), its edge valence is
\[
\deg_K(e)=\#\{\sigma\in K:\dim\sigma=3,\ e\subseteq\sigma\}.
\]
If a triangulation instead permits identifications among faces of abstract tetrahedra, edge valence counts incident tetrahedral occurrences with multiplicity. The supplied record does not specify which triangulation convention is intended.

The source asks whether every orientable \(3\)-manifold \(M\) admits a triangulation \(T\) satisfying
\[
\forall e\in T^{(1)}_{\mathrm{edges}},\qquad \deg_T(e)\in\{5,6\},
\]
where \(T^{(1)}_{\mathrm{edges}}\) denotes the set of edges. This condition allows either valence to be absent; it does not require both valences to occur.

There is an unresolved scope qualification: the extracted statement says every orientable \(3\)-manifold, without specifying compactness or boundary, whereas the supplied literature assessment restricts the question to closed orientable \(3\)-manifolds. Here closed means compact and without boundary. Under the ordinary simplicial convention, that restricted formulation is
\[
\forall M\text{ a closed orientable }3\text{-manifold},\quad
\exists K\text{ a finite simplicial complex},\quad
|K|\cong M\quad\text{and}\quad
\forall e\in K,\ \dim e=1\Longrightarrow\deg_K(e)\in\{5,6\}.
\]
This is an explicitly qualified formulation, not a resolution of the discrepancy with the broader extracted statement. Connectedness, the treatment of manifolds with boundary, and local finiteness for any noncompact version are not specified in the record and require review.

The supplied background describes an established construction with edge valences in \(\{4,5,6\}\), and presents the requested restriction to \(\{5,6\}\) as an open TCP-triangulation problem for closed orientable \(3\)-manifolds. These are supplied contextual and status claims, not independently verified mathematical or literature conclusions.
\subsection{Short English statement}
Can every orientable \(3\)-manifold be triangulated so that each edge has valence \(5\) or \(6\)? The supplied literature assessment asks this for closed orientable manifolds and describes a known \(4,5,6\) construction. The intended manifold scope and triangulation convention remain unresolved.
\subsection{Sources}
\begin{itemize}
\item Supplied local UnsolvedMath registry record, permanent ID \(30002070\), source problem code \texttt{catalogue entry}, in the Oberwolfach Reports Open Problems collection. The record labels the statement exact and reports manual review of the extraction; these are supplied verification annotations. No individual proposer is identified.
\item Supplied report citation: \emph{Triangulations (2013)}, Oberwolfach Reports, pp.\ 1405--1486, DOI \texttt{10.4171/owr/2012/24}. \url{https://doi.org/10.4171/owr/2012/24}. The registry records proposed year \(2012\); the supplied citation retains \(2013\) in its title.
\item Supplied literature source: Brady--McCammond--Meier. \url{https://doi.org/10.2307/1194005}. No fuller bibliographic description is supplied.
\item Supplied literature source labelled ``original OWR report'': \url{https://doi.org/10.4171/owr/2012/24}.
\item Literature assessment dated exactly \texttt{2026-08-21}: the record labels the problem open, describes the \(4\)--\(5\)--\(6\) edge-degree theorem as established, and describes the \(5\)--\(6\) assertion for every closed orientable \(3\)-manifold as an open TCP-triangulation problem. This is supplied dated curation derived from the OWR material, not an independently verified assessment. The supplied background explicitly calls for independent verification of time-sensitive status and cited claims.
\end{itemize}
\end{document}
